\chapter{总结与展望}
\label{chap:6}

\section{论文总结}

本文面向 Linux 内核 built-in 代码这一被现有部分随机化方案在原理上遗漏的部分，提出并实现了一套持续随机化系统。第 1.2 节以两条独立证据说明了这一遗漏的严重性：从源码层面看，内核最核心、执行最频繁的子系统在任何配置下都无法编译为模块，完全落在以模块为对象的方案之外；从漏洞分布看，2024 年 1 月至 2026 年 9 月公开的 10,705 个内核 CVE 中有 21.6\% 涉及只能内建的代码。built-in 代码没有模块那样天然的随机化边界，其代码索引更新的规模、执行流的安全暂停、以及随机化范围之外代码的保护，都缺乏可用的现成解法。本文以跳板函数作为随机化区域的唯一出入口，在其之上组织三项关键技术，分别回答持续随机化的三个核心问题。

\textbf{第一项工作回答“如何随机化”。} 第 3 章提出基于跳板函数的低开销持续随机化执行机制。它在保持内核原有代码模型不变的前提下，由编译器插件按函数改写被选中函数内部的引用，使函数体可以整体搬移；跳板在承担地址间接的同时充当执行流的观测点，由进出计数精确刻画随机化区域内的活跃执行流，从而实现不保留旧副本的切换。这一设计把一次随机化需改写的代码索引由与调用点数量同阶降到与函数数量同阶，实测一次跨区域调用往返的常驻开销为 18.02 ns，其中执行流追踪占绝大部分而位置无关改造本身仅占 0.5 至 1.4 ns。第 3 章同时给出了这条路线的边界：全树 9.6\% 的函数因 per-CPU 变量的访问方式而无法随机化，变参函数与静态调用点另有各自的不相容之处，本章机制与间接分支防护类机制之间还存在结构性互斥。

\textbf{第二项工作回答“何时随机化”。} 第 4 章把只执行内存由阻断机制改用为检测机制，以对代码探测行为的直接观测替代时钟或行为代理指标作为触发条件，并给出基于管理态保护键与基于硬件调试观察点的两条实现路径，两条路径的单次检测开销分别为 0.96 μs 与 1.32 μs。这一改变不仅提供了精准的触发时机，还化解了纯防护式只执行内存在内核中最棘手的问题：合法的代码读取路径不再需要被逐一豁免，漏掉一条只是多触发一次随机化，而不再是内核崩溃。正常负载下的触发频率为每秒 0.02 至 0.44 次，从检测到随机化完成的延迟在基础负载下 99 分位不超过 110 μs；竞态实验表明，当读取至解引用的时延达到 92 μs 时，攻击成功率已降至 1.0\%。

\textbf{第三项工作回答范围之外代码的保护。} 第 5 章利用随机化区域内目标地址不可预测这一性质，提出以随机化打断指针认证上下文传播链的协同机制：以类型上下文恒定签名配合归属跟踪在运行时识别高扇入函数并将其移入随机化区域，并在跳板处施加分区验证，进入随机化区域免验证、离开时强制认证与白名单检查。实测 91\% 的间接调用点合法目标数不超过 5，最大等价类由类型哈希方案的 141 降到 57；一次跨区域离开的开销为 9.0 ns。

\textbf{系统层面}，三项机制共享同一套基础设施：跳板运行时、变体池与布局分配器、跨区域白名单。第 4 章的触发式随机化建立在第 3 章的跳板之上，只需在一级跳板处改写一条分支指令即可切换激活状态；第 5 章的分区验证同样施加在第 3 章已有的跨区域出入口上，不另建检查点。这一共享使后两章叠加在第 3 章之上的边际开销很小，也是三项工作能够在同一平台上构成一个完整可运行系统的原因。

\section{主要创新点}

本文的贡献可归纳为三项机制创新与一项系统创新。

\textbf{其一，把跳板从单纯的地址间接层扩展为执行状态的观测点。} 已有工作使用跳板是为了间接寻址，而本文让同一份基础设施在承担间接寻址的同时刻画随机化区域内的活跃执行流，由此得到不保留旧副本的切换，避免了旧副本在被回收前作为攻击者可用代码来源的暴露窗口。与之配套，本文在不改变内核代码模型的前提下，用编译器插件按函数施加位置无关改造，解决了 built-in 代码没有模块边界、无法借助加载器重定位这一根本困难。

\textbf{其二，把只执行内存从阻断改用为检测，使随机化的触发条件成为对攻击行为的直接观测。} 这把“何时随机化”从依赖时钟或行为代理指标，转变为依赖“代码是否真的被探测过”这一可观测事实，在无攻击时几乎不付出代价、攻击一旦发生则立即响应。两条基于不同硬件原理的实现路径说明该机制不绑定特定架构，其依赖的只是“能够观测代码读取”这一能力。

\textbf{其三，把随机化与控制流完整性从两个独立叠加的机制变为互为条件的协同机制。} 随机化区域内目标地址不可预测，因而进入该区域的转移不需要验证；把高扇入函数移入随机化区域，即可从中间打断指针认证的上下文传播链，使两侧调用点各自建立唯一上下文。随机化为控制流完整性提供了打断传播链的手段，控制流完整性为随机化补上了范围之外的防护。

\textbf{在系统层面}，三项工作不是三个独立系统的拼装，而是共享同一套基础设施（跳板运行时、变体池与布局分配器、跨区域白名单）的三层机制。这一共享使第 4、5 两章叠加在第 3 章之上的边际开销很小，也是三项工作能够在同一平台上构成一个完整可运行系统的原因。

\section{局限性与未来工作}

本文在各章分别声明了机制的边界，此处汇总为四类，以便完整评估本文工作的适用范围。

\textbf{能力边界：存在无法随机化的函数。} 全树 9.6\% 的函数因 per-CPU 变量必须以常量偏移访问而无法被随机化，这一比例是方法的固有边界而非实现未完成；变参函数（150 个）受跳板转发能力限制，静态调用点在使用它们的配置下与位置无关性本质冲突（在本文评测配置下为零）。这些函数留在非随机化区域，其防护由第 4、5 章承担，而非靠“把所有代码都随机化”。

\textbf{代价边界：常驻开销随范围上升。} 常驻开销主要由跨区域调用的频度决定，随随机化范围扩大而上升；把范围扩到全部内建代码将不再具备部署价值，因此范围必须限定。这与能力边界共同决定了非随机化区域的存在是方法本身的结果，而非取舍。

\textbf{功能边界：动态跟踪失效与间接分支防护互斥。} 被随机化的函数上，动态跟踪类设施（ftrace、kprobes）表现为静默失效而非报错；原因是这些设施读取并改写函数入口的原指令，而被随机化的函数已整体搬离该位置。本文机制为实现位置无关而制造间接调用，与 x86 的间接分支防护、arm64 的分支目标识别存在结构性互斥，评测中一律关闭且四档保持一致。此外，第 4 章改写跳转桩时经 \lstinline[breaklines=false]@text_poke@ 建立的临时可写映射，在其存在期间是跳转桩区的第三份别名，构成一个每核可见、关中断期间存在的短暂窗口。

\textbf{假设边界：指针认证的前提是脆弱的。} 第 5 章的完整性论证以“攻击者无法获得验证预言机”为前提，而推测执行侧信道类攻击（PACMAN\citep{ravichandran2022pacman}）恰恰提供了这样一个预言机；本文不声称在此类攻击存在时不变式 I3 仍然成立，相关防御与本文正交。运行时识别高扇入函数只能覆盖实际执行到的路径，冷路径上的高扇入函数不会被发现，其调用点长期停留在共享的类型上下文上，这一漏识别在本机制中同时构成可用性风险。

上述四条边界各自指向一个可以继续推进的方向，以下四项按优先次序给出。

\textbf{其一，缩小无法随机化的比例。} per-CPU 访问是当前唯一大量存在的排除来源。逐条改写内核的 per-CPU 访问器宏，使其偏移可被物化，理论上可解，但代价与收益是否相称需要评估；一个更现实的方向是把非随机化区域内的这部分函数交由更强的控制流完整性检查覆盖，而非强行随机化。

\textbf{其二，降低常驻开销的主项。} 第 3 章的开销归因表明执行流进出计数一项独占常驻开销的近七成，当前以一个全局原子量实现；改为每处理器核独立计数、由随机化线程求和，是第一顺位的优化。指针认证的上下文计算可由编译期常量折叠进一步降低。二者若均落实，端到端开销有望再降约一半。

\textbf{其三，把两条硬件路径推向更广的机型。} 两条路径均已在各自平台上实现并测得单次检测开销，但观察点在不同处理器实现上的行为差异、保护键在不同微架构上的写寄存器代价都还需要更广的验证。把两条路径在同一套负载上做跨机型的完整对照，是使第 4 章结论更可部署的必要一步。

\textbf{其四，正面应对推测执行侧信道。} 指针认证在推测执行下可被爆破这一前提，是第 5 章论证中最脆弱的一环。把针对推测执行的既有防御与本文机制组合部署、并评估其叠加代价，是使完整性不变式在更强攻击者模型下依然成立的方向。

\section{结语}

本文的出发点是一个具体而被长期回避的问题：内核中最核心、执行最频繁、也最常出漏洞的那部分代码，恰恰落在现有随机化方案的保护范围之外。本文没有试图用一种机制解决全部困难，而是承认 built-in 代码必然留下无法随机化的部分，并围绕这一事实组织出一条约束链：因为存在无法随机化的函数、且把范围扩到全部代码的开销不可接受，所以非随机化区域必然存在；于是随机化区域的地址如何不被泄露、何时触发才既及时又不浪费，以及非随机化区域内已知布局的代码如何不被复用，被依次逼出，分别由本文的三项工作回答。三者共享同一套跳板基础设施，因而不是三种防护的并列叠加，而是同一条约束链上被依次逼出的三步。本文期望这一以承认边界为起点、以协同而非叠加为组织方式的思路，能为 built-in 内核代码的持续保护提供一个可部署、可复现、边界清晰的参照。

